
\SGARefTarget{sga3:xi:section:3}\section*{3. Auxiliary Representability
Results}
\addcontentsline{toc}{section}{3. Auxiliary Representability Results}

\medskip
\noindent\SGARefTarget{sga3:xi:proposition:3:1}\textbf{Proposition 3.1.} ---
\textit{Let \(F\to S\) be a functor over the prescheme \(S\).  The
following conditions are equivalent:}

\begin{enumerate}
\item[(i)] \SGARefTarget{sga3:xi:proposition:3:1:item:i}
\textit{\(F\) is representable and \(F\to S\) is an open immersion
(one also says simply that \(F\to S\) is an open immersion).}

\item[(ii)] \SGARefTarget{sga3:xi:proposition:3:1:item:ii}
\textit{\(F\) is a sheaf for the faithfully flat quasi-compact
topology, ``commutes with filtered direct limits of rings,'' and
\(F\to S\) is a monomorphism.  Moreover, the following condition
holds: for every local prescheme \(S'\) over \(S\), with residue field
\(k\), and every \(S\)-morphism
\[
  \operatorname{Spec}(k)=S'_0\longrightarrow F,
\]
there exists an \(S\)-morphism \(S'\to F\) extending it.\footnote{The
French re-edition prints \(S'\to S\).  Since this morphism must extend
the displayed morphism \(S'_0\to F\), its intended target is \(F\).}}

\item[(iii)] \SGARefTarget{sga3:xi:proposition:3:1:item:iii}
\textit{Suppose that \(S\) is locally noetherian.  Then \(F\) is a
sheaf for the faithfully flat quasi-compact topology, commutes with
filtered direct limits of rings, commutes with adic inverse limits of
local artinian rings, and \(F\to S\) is a monomorphism and is
infinitesimally étale
(cf.\ \SGARefLink{sga3:xi:remark:1:8}{1.8}).}
\end{enumerate}

Let us first make two points of terminology precise.

\medskip
\noindent\SGARefTarget{sga3:xi:remark:3:2}\textbf{Remark 3.2.} ---
A functor \(F\) over \(S\) is said to \emph{commute with
direct limits of rings} (understood to mean filtered direct limits) if
the following holds.  Let \((S'_i)_{i\in I}\) be any filtered inverse
system of affine \(S\)-schemes lying over an affine open subscheme of
\(S\), and let \(A'_i\) be their rings.  Then the natural homomorphism

\SGARefTarget{sga3:xi:display:3:2:direct-limit}
\[
  \tag{*}
  \varinjlim_i F(S'_i)\longrightarrow F(S'),
  \qquad
  S'=\operatorname{Spec}(A'),
  \quad
  A'=\varinjlim_i A'_i
\]
is bijective.

Notice that \(S'\) is precisely the inverse limit of the \(S'_i\) in
the category of preschemes, and even in the category of all ringed
spaces.  Thus the condition under consideration has the character of
a right-exactness condition---commutation with certain direct limits
in \((\mathbf{Sch})^\circ/S\)---just as does the condition of being a
sheaf for a topology.

Care must be taken that this condition is essentially relative: it
involves the morphism \(F\to S\), not merely the functor
\[
  F:(\mathbf{Sch})^\circ\longrightarrow(\mathbf{Ens}).
\]
More precisely, in \SGARefLink{sga3:xi:display:3:2:direct-limit}{\((*)\)}, \(F(S'_i)\) and \(F(S')\) mean
\(\operatorname{Hom}_S(S'_i,F)\) and
\(\operatorname{Hom}_S(S',F)\).  Consequently, when \(F\) is
representable, the condition says that \(F\) is locally of finite
presentation over \(S\).  We used several times in the preceding two
exposés the fact that a functor represented by an \(S\)-prescheme
locally of finite presentation commutes with direct limits of rings.

\medskip
\noindent\SGARefTarget{sga3:xi:remark:3:3}\textbf{Remark 3.3.} ---
A functor \(F\) over \(S\) is said to \emph{commute with adic inverse
limits of local artinian rings} if the following condition holds.
For every \(S'\) over \(S\) that is the spectrum of a complete
noetherian local ring \(A'\), set
\[
  S'_n=\operatorname{Spec}
  \bigl(A'/\operatorname{rad}(A')^{\,n+1}\bigr).
\]
Then the natural map

\SGARefTarget{sga3:xi:display:3:3:adic-limit}
\[
  \tag{**}
  F(S')\longrightarrow\varprojlim_n F(S'_n)
\]
is required to be bijective.

This condition has the character of a left-exactness condition and
is satisfied whenever \(F\) is representable.  In contrast to
commutation with direct limits of rings, it is intrinsic to \(F\) as
an object of \(\widehat{\mathbf{Sch}}\); that is, it does not involve
the morphism \(F\to S\).

\medskip
\noindent\SGARefTarget{sga3:xi:remark:3:4}\textbf{Remark 3.4.} ---
Let \(F\) be a functor over \(S\) that is a sheaf for the Zariski
topology, or, as one also says, is ``of local nature.''  It is enough
for this that \(F\) be a sheaf for a finer topology, such as the
faithfully flat quasi-compact topology.  Let \((S_i)\) be an open
covering of \(S\).  A straightforward gluing argument shows that
\(F\) is representable if and only if each
\[
  F_i=F\times_S S_i
\]
is representable; this allows one, for example, to reduce to the case
where \(S\) is affine.

Suppose that the functor \(F\), of local nature, commutes with direct
limits of rings.  Then \(F\) is representable if and only if its
restriction to the category of preschemes locally of finite
presentation over \(S\) is representable.  Necessity was noted in
\SGARefLink{sga3:xi:remark:3:2}{3.2}.  Sufficiency amounts to the
following assertion: if \(X\) is a prescheme locally of finite
presentation over \(S\), and \(X\to F\) is a morphism such that, for
every \(S'\) locally of finite presentation over \(S\), the induced
map

\SGARefTarget{sga3:xi:display:3:4:hom-comparison}
\[
  \operatorname{Hom}_S(S',X)\longrightarrow
  \operatorname{Hom}_S(S',F)
\]
is bijective, then \(X\to F\) is an isomorphism.  This follows readily
from the fact that \(X\) and \(F\) are both functors of local nature
that commute with direct limits of rings.

\medskip
\noindent\textit{Proof of \SGARefLink{sga3:xi:proposition:3:1}{Proposition 3.1}.}
The implications \SGARefLink{sga3:xi:proposition:3:1:item:i}{(i)} \(\Rightarrow\) \SGARefLink{sga3:xi:proposition:3:1:item:ii}{(ii)} and
\SGARefLink{sga3:xi:proposition:3:1:item:i}{(i)} \(\Rightarrow\) \SGARefLink{sga3:xi:proposition:3:1:item:iii}{(iii)} are immediate.  We prove the converse
implications.

Assume \SGARefLink{sga3:xi:proposition:3:1:item:ii}{(ii)}.  Let \(U\) be the set of points \(s\in S\) for which the
canonical monomorphism
\[
  \operatorname{Spec}\kappa(s)\longrightarrow S
\]
factors through \(F\).  By the last condition in \SGARefLink{sga3:xi:proposition:3:1:item:ii}{(ii)}, for every
\(s\in U\), the canonical map
\(\operatorname{Spec}\mathscr O_{S,s}\to S\), equivalently its
factorization through \(U\), factors through \(F\).  The local ring
\(\mathscr O_{S,s}\) is the direct limit of the rings of affine
neighborhoods of \(s\).  Since \(F\) commutes with direct limits of
rings, for every \(s\in U\)\footnote{The French re-edition prints
\(s\in S\) at this occurrence; the argument and the conclusion
\(U^s\subset U\) require \(s\in U\).} there is an open neighborhood
\(U^s\) such that the canonical immersion \(U^s\to S\) factors
through \(F\).  It follows that \(U^s\subset U\), and hence that
\(U\) is open.

Because \(F\to S\) is a monomorphism and \(F\) is of local nature,
the \(S\)-morphisms \(U^s\to F\) agree on the intersections
\(U^s\cap U^{s'}\), for \(s,s'\in U\), and therefore glue to an
\(S\)-morphism \(U\to F\).  It remains to prove that this morphism is
an isomorphism, or equivalently that every \(S\)-morphism \(S'\to F\)
factors uniquely through \(U\), where \(S'\) is an \(S\)-prescheme.
Since both \(F\to S\) and \(U\to S\) are monomorphisms, it is
equivalent to say that the structure morphism \(S'\to S\) factors
through \(U\).

By locality, we are reduced to the case in which \(S'\) is the
spectrum of a field and hence has a single point \(s'\).  Let \(s\)
be the point of \(S\) below \(s'\).  We claim that the given
\(S\)-morphism \(S'\to F\) factors through
\[
  \operatorname{Spec}\kappa(s)=S_0\longrightarrow F.
\]
This implies \(s\in U\) and proves the assertion.  Since
\(S'\to S_0\) is an fpqc covering and \(F\) is a sheaf for the fpqc
topology, it is enough that the two composites

\SGARefTarget{sga3:xi:display:3:1:fpqc-descent}
\[
  S''=S'\times_{S_0}S'
  \rightrightarrows S'\longrightarrow F
\]
coincide.  They do coincide because \(F\to S\) is a monomorphism.
This proves \SGARefLink{sga3:xi:proposition:3:1:item:ii}{(ii)} \(\Rightarrow\) \SGARefLink{sga3:xi:proposition:3:1:item:i}{(i)}.

Now assume \SGARefLink{sga3:xi:proposition:3:1:item:iii}{(iii)}, with \(S\) locally noetherian.  It is enough to
prove the last condition of \SGARefLink{sga3:xi:proposition:3:1:item:ii}{(ii)}; by the preceding proof, it is
enough to do so when
\[
  S'=\operatorname{Spec}\mathscr O_{S,s}
\]
for some \(s\in S\).  Put
\[
  A=\mathscr O_{S,s},
  \qquad
  A_n=A/\mathfrak m^{\,n+1},
  \qquad
  S_n=\operatorname{Spec}(A_n).
\]
Because \(F\to S\) is infinitesimally smooth, the given morphism
\(S'_0\to F\) extends to morphisms \(S_n\to F\).  Since \(F\to S\)
is a monomorphism, these morphisms form an element of
\(\varprojlim_n F(S_n)\).  Since \(F\) commutes with adic inverse
limits of local artinian rings, they arise from a morphism
\[
  \operatorname{Spec}(\widehat A)=\widehat S'\longrightarrow F.
\]
Finally, \(F\to S\) is a monomorphism, \(F\) is an fpqc sheaf, and
\(\widehat S'\to S'\) is an fpqc covering.  The last morphism
therefore descends to \(S'\to F\), completing the proof of
\SGARefLink{sga3:xi:proposition:3:1:item:iii}{(iii)} \(\Rightarrow\) \SGARefLink{sga3:xi:proposition:3:1:item:ii}{(ii)}. \hfill\(\square\)

\medskip
\noindent\SGARefTarget{sga3:xi:proposition:3:5}\textbf{Proposition 3.5.} ---
\textit{Let \(S\) be a locally noetherian prescheme, let \(F\to S\)
be a functor over \(S\), and let
\((X_i,u_i)_{i\in I}\) be a family of \(S\)-morphisms
\(u_i:X_i\to F\), where the \(X_i\) are preschemes locally of finite
type over \(S\).  Suppose that the following conditions hold:}

\begin{enumerate}
\item[(a)] \SGARefTarget{sga3:xi:proposition:3:5:item:a}
\textit{\(F\) is a sheaf for the faithfully flat quasi-compact
topology, commutes with direct limits of rings, and commutes with adic
inverse limits of local artinian rings.}

\item[(b)] \SGARefTarget{sga3:xi:proposition:3:5:item:b}
\textit{The morphisms \(u_i:X_i\to F\) are monomorphisms and are
infinitesimally étale
(cf.\ \SGARefLink{sga3:xi:remark:1:8}{1.8}).}

\item[(c)] \SGARefTarget{sga3:xi:proposition:3:5:item:c}
\textit{The family \((u_i)\) is ``set-theoretically surjective.''}
\end{enumerate}

\textit{Under these conditions, \(F\) is representable by a
prescheme locally of finite type over \(S\).  Moreover, the \(u_i\)
are open immersions and the family \((X_i)\) is an open covering of
\(F\).}

\medskip
\noindent\SGARefTarget{sga3:xi:remark:3:6}\textbf{Remark 3.6.} ---
Proceeding as at the end of
\SGARefLink{sga3:xi:remark:1:7}{Remark 1.7}, one defines, for every
functor
\[
  F:(\mathbf{Sch})^\circ\longrightarrow(\mathbf{Ens}),
\]
an ``underlying set'' \(\operatorname{ens}(F)\): it is the quotient
of the set of field-valued points of \(F\) by the equivalence relation
described in \SGARefLink{sga3:xi:remark:1:7}{1.7}.  When \(F\) is represented by \(X\), this recovers
the underlying set of \(X\).

Plainly, \(\operatorname{ens}(F)\) depends functorially on \(F\).
Thus, if \(G\to F\) is a morphism of functors, we say that it is
\emph{set-theoretically surjective} if the induced map
\[
  \operatorname{ens}(G)\longrightarrow\operatorname{ens}(F)
\]
is surjective.  Equivalently, every point of \(F\) with values in a
field \(k\) ``comes from'' a point of \(G\) with values in a suitable
extension of \(k\).  This definition extends immediately to a family
of morphisms \(G_i\to F\), and thereby specifies the meaning of
condition \SGARefLink{sga3:xi:proposition:3:5:item:c}{(c)}.

\medskip
\noindent\textit{Proof of \SGARefLink{sga3:xi:proposition:3:5}{Proposition 3.5}.}
For \((i,j)\in I\times I\), put

\SGARefTarget{sga3:xi:display:3:5:fiber-product}
\[
  X_{i,j}=X_i\times_F X_j
\]
and consider the projections

\SGARefTarget{sga3:xi:display:3:5:projections}
\[
  v_{i,j}:X_{i,j}\longrightarrow X_i,
  \qquad
  w_{i,j}:X_{i,j}\longrightarrow X_j.
\]
We claim that these morphisms are representable by open immersions.
Apply criterion
\SGARefLink{sga3:xi:proposition:3:1:item:iii}{3.1(iii)}.
The functor \(X_{i,j}\) satisfies the three exactness conditions:
it is a sheaf, it commutes with direct limits of rings, and it
commutes with adic inverse limits of local artinian rings.  Indeed,
\(F\), \(X_i\), and \(X_j\) satisfy these conditions, and the
conditions are stable under finite inverse limits, in particular
under fiber products.  Since every \(X_i\to F\) is a monomorphism,
both projections are monomorphisms, each being obtained from one of
these maps by base change.  Finally, being infinitesimally étale is
also preserved by base change.  Thus the hypotheses of \SGARefLink{sga3:xi:proposition:3:1:item:iii}{3.1(iii)} are
satisfied.

We may now use the \(X_i\), \(X_{i,j}\), \(v_{i,j}\), and \(w_{i,j}\)
to construct in the usual way an \(S\)-prescheme \(X\), in which the
\(X_i\) are identified with open subschemes, the \(X_{i,j}\) with the
intersections \(X_i\cap X_j\), and \(v_{i,j}\), \(w_{i,j}\) with the
canonical immersions.

Put\footnote{The French re-edition labels this coproduct \(X\) and
later writes \(R=X\times_F X\), while simultaneously using
\(R\rightrightarrows Y\) and \(X=Y/R\).  The coherent notation,
adopted here, is \(Y=\coprod_iX_i\) and \(R=Y\times_FY\).}

\SGARefTarget{sga3:xi:display:3:5:coproduct}
\[
  Y=\coprod_i X_i.
\]
Then \(X\) is also the quotient of \(Y\) by the equivalence relation

\SGARefTarget{sga3:xi:display:3:5:relation}
\[
  R=\coprod_{i,j}X_{i,j}
  \rightrightarrows Y,
\]
whose two projections \(v,w\) are defined by the \(v_{i,j}\) and
\(w_{i,j}\), respectively.

More precisely, because \(F\) is an fpqc sheaf, the morphisms
\(u_i:X_i\to F\) induce a morphism \(u:Y\to F\), and

\SGARefTarget{sga3:xi:display:3:5:kernel-relation}
\[
  R=Y\times_FY
\]
is exactly the equivalence relation defined by \(u\).  The quotient
\(X=Y/R\) is also a quotient in the category of fpqc sheaves, and
even in the category of Zariski sheaves; this follows directly from
the definitions of quotient and sheaf.  Consequently, \(u\) factors
uniquely through a morphism

\SGARefTarget{sga3:xi:display:3:5:induced-map}
\[
  \overline u:X\longrightarrow F,
\]
and \SGARefLink{sga3:xi:display:3:5:induced-map}{this morphism} is a monomorphism.

It remains to show that \(\overline u\) is an isomorphism.  Since
\(F\) is of local nature, we may assume that \(S\) is affine.  Since
\(F\) also commutes with direct limits of rings, it is enough, by
\SGARefLink{sga3:xi:remark:3:4}{3.4}, to verify that, for every
affine \(T\) of finite type over \(S\), every morphism \(T\to F\)
factors through \(X\).  Consider

\SGARefTarget{sga3:xi:display:3:5:base-change}
\[
  G=X\times_FT\longrightarrow T.
\]
It is enough to prove that \SGARefLink{sga3:xi:display:3:5:base-change}{this morphism} is an isomorphism.  Since
\(T\) is noetherian, the argument above shows that it is an open
immersion.  Here one uses that \(X\to F\) is infinitesimally étale,
which follows immediately because the induced morphisms
\(u_i:X_i\to F\) are infinitesimally étale.  By hypothesis,
\(X\to F\) is set-theoretically surjective, and this condition is
stable under base change.  Hence \(G\to T\) is set-theoretically
surjective.  Being also an open immersion, it is an isomorphism.
\hfill\(\square\)

